% homework/alg13/alg13.tex
% Algebra Track 1.3 - The Circle Map.
%
% NOT a Year 7 packet. This is a ONE-PAGE, LANDSCAPE handout for the high-school
% Algebra Track, used in class during lesson 1.3 and printed from the site's
% Homework shelf. It has no sections, no marks, no answer lines to speak of --
% the whole page exists to give a student a map with room to swing a compass on.
%
% Landscape is not decoration. The sheet is worked on with a compass, and the
% arcs of a perpendicular-bisector construction run well outside the circle they
% produce, so the binding constraint is drawing room. Landscape puts the map at
% 18.6 cm wide against 17 cm in portrait, and moves the instructions into a side
% column instead of stealing height from the map.
%
% The two tasks mirror the lesson deliberately. Task 1 is done BY EYE, before the
% construction has been taught (deck slides 6-7): everyone succeeds roughly and
% nobody succeeds exactly, which is what makes the compass worth learning. Task 2
% is the same sheet after the construction (deck slide 23), and the comparison
% between the two circles is the assessment.
%
% THE CITY COORDINATES ARE NOT EYEBALLED. europe-map.png is an equirectangular
% crop at a uniform 30 px/degree covering 12W-30E and 36N-64N, so it is 1260x840
% and a city sits at x = (longitude + 12) * 30, y = (64 - latitude) * 30. The
% \city macro takes those figures normalised to the 0..1 box that TikZ overlays
% on the image, i.e. (x/1260, 1 - y/840). Same numbers as the deck's widget --
% see content/freshman-math/U01_3/images/CREDITS.json.
%
% The source is pure ASCII: every typographic character goes through a LaTeX
% command, so no editor or shell can mangle the encoding.
%
% Build:  pdflatex alg13.tex   (run twice, for the TikZ bounding box)

\documentclass[12pt,a4paper]{article}

% -- Answer key switch -------------------------------------------------------
% There is no answer key for this sheet -- there is no single right answer, and
% the teacher notes live in the lesson plan. The switch is kept so the shared
% build script can run its teacher pass without failing.
\newif\ifkey
\ifdefined\TEACHER\keytrue\else
  \keyfalse
\fi

% -- House style -------------------------------------------------------------
% homework/hw-style.tex
% Shared house style for every Year 7 homework packet. \input this from the
% packet file, right after \documentclass and the answer-key switch.
%
% Do not fork this file per packet. If a packet needs a different look, the
% question is whether every packet should have it.
%
% The choices here are deliberate and were arrived at by printing and looking:
%   * No solid dark banners. Every header is a light tint with a coloured spine
%     on the left, so colour reads clearly without flooding a school printer.
%   * No marks and no mark totals. These are practice, not exam papers.
%   * Lato at 12pt with open leading, plus mathastext so inline maths is Lato
%     too. Without mathastext every "-6 + 4" is Computer Modern serif and the
%     sheet looks half-typeset.
%   * Writing lines are spaced for Year 7 handwriting, not adult margin notes.
%
% Provides: \wlines \blank \tickbox \sep \qmk-free marks (none) \hwsection
% \qhead \expr, the workspace box, and the remember / wordhelp / activity /
% yourturn coloured boxes. Palette matches the lesson decks exactly.

% -- Packages ----------------------------------------------------------------
\usepackage[T1]{fontenc}
\usepackage[utf8]{inputenc}
\usepackage[default]{lato}
\usepackage[a4paper,top=1.7cm,bottom=1.7cm,left=2.0cm,right=2.0cm,headheight=15pt]{geometry}
\usepackage{amsmath,amssymb}
\usepackage{xcolor}
\usepackage[most]{tcolorbox}
\usepackage{enumitem}
\usepackage{tikz}
\usetikzlibrary{arrows.meta}
\usepackage{array}
\usepackage{tabularx}
\usepackage{fancyhdr}
\usepackage{multicol}
\usepackage{needspace}
\usepackage{graphicx}

% Make maths use Lato too. Without this, every "-6 + 4" on the page is set in
% Computer Modern serif and the sheet looks half-typeset. Must come last.
\usepackage[italic]{mathastext}

\linespread{1.06}\selectfont

% -- The Learner's Book palette (same hex values as the lesson decks) ---------
\definecolor{cbteal}{HTML}{0087A8}
\definecolor{cbpurple}{HTML}{5C2483}
\definecolor{cborange}{HTML}{C25E12}
\definecolor{cbgreen}{HTML}{4A8B23}
\definecolor{cbred}{HTML}{C8102E}
\definecolor{cbblue}{HTML}{1A5FA8}
\definecolor{rulegray}{HTML}{9AA5AD}

% -- Answer lines and blanks -------------------------------------------------
% \wlines{n} draws n ruled writing lines, spaced wide enough for Year 7
% handwriting rather than for an adult with a fine pen.
\makeatletter
\newcommand{\wlines}[1]{%
  \par\vspace{0.75em}%
  \@tempcnta=0\relax
  \@whilenum\@tempcnta<#1\do{%
    \hbox to \linewidth{\color{rulegray}\leaders\hrule height 0.5pt\hfill}%
    \vspace{1.5em}%
    \advance\@tempcnta by 1\relax}%
  \par\vspace{-0.8em}%
}
\makeatother

% \blank[width] is an inline gap to write a short answer in.
\newcommand{\blank}[1][2.6cm]{\underline{\hspace{#1}}}

% A boxed space for showing working.
\newtcolorbox{workspace}[1][3.4cm]{%
  colback=white, colframe=rulegray, boxrule=0.5pt, arc=5pt,
  height=#1, left=6pt, right=6pt, top=4pt, bottom=4pt,
  before skip=9pt, after skip=11pt}

% An empty tick box.
\newcommand{\tickbox}{\fbox{\rule{0pt}{1.15em}\rule{1.15em}{0pt}}}

% A middot separator.
\newcommand{\sep}{\,\textperiodcentered\,}

% -- Section and question headings -------------------------------------------
% Light tint plus a fat coloured spine on the left. No solid dark banner: the
% colour reads clearly but barely costs the printer anything.
% needspace stops a heading stranding alone at the foot of a page. Kept modest:
% too large and it pushes headings over, leaving a worse hole than it prevents.
\newcommand{\hwsection}[2][cbteal]{%
  \needspace{8\baselineskip}%
  \par\vspace{13pt}%
  \noindent
  \begin{tcolorbox}[enhanced, nobeforeafter, width=\textwidth,
      colback=#1!7, colframe=#1!7, arc=5pt,
      borderline west={5pt}{0pt}{#1},
      left=14pt, right=10pt, top=7pt, bottom=7pt]
    {\large\bfseries\color{#1}#2}
  \end{tcolorbox}%
  \par\vspace{9pt}}

% \qhead[colour]{A2}{Moving along the line} -- a question heading.
\newcommand{\qhead}[3][cbteal]{%
  \needspace{6\baselineskip}%
  \par\vspace{11pt}%
  \noindent{\large\bfseries\color{#1}#2}\quad{\large\bfseries #3}\par\vspace{4pt}}

% -- Coloured boxes ----------------------------------------------------------
% Shared look: rounded, light fill, coloured rule, coloured title text on a
% pale title strip rather than a dark one.
\tcbset{hwbox/.style={
  enhanced, breakable, arc=6pt, boxrule=1pt,
  left=11pt, right=11pt, top=8pt, bottom=8pt,
  fonttitle=\bfseries, attach boxed title to top left={xshift=12pt, yshift=-8pt},
  boxed title style={arc=4pt, boxrule=1pt, left=7pt, right=7pt, top=2pt, bottom=2pt},
  top=13pt}}

% Orange: things to remember. Matches the deck's "Write This Down".
\newtcolorbox{remember}[1][Remember from the lesson]{hwbox,
  colback=cborange!6, colframe=cborange, coltitle=cborange,
  colbacktitle=white, title={#1}}

% Blue: English help. The barrier in this class is the wording, not the maths.
\newtcolorbox{wordhelp}[1][Word help]{hwbox,
  colback=cbblue!6, colframe=cbblue, coltitle=cbblue,
  colbacktitle=white, title={#1}}

% Purple: an activity or instruction.
\newtcolorbox{activity}[1][How to do this homework]{hwbox,
  colback=cbpurple!6, colframe=cbpurple, coltitle=cbpurple,
  colbacktitle=white, title={#1}}

% Green: your turn / a friendly nudge.
\newtcolorbox{yourturn}[1][Your turn]{hwbox,
  colback=cbgreen!7, colframe=cbgreen, coltitle=cbgreen!75!black,
  colbacktitle=white, title={#1}}

% -- Shorthands --------------------------------------------------------------
\newcommand{\dC}{\,^{\circ}\mathrm{C}}          % use inside maths: $-8\dC$
\newcommand{\key}[1]{{\color{cborange}\textbf{#1}}}
\newcommand{\eng}[1]{{\color{cbblue}\textbf{#1}}}

\setlist[enumerate]{itemsep=8pt, topsep=5pt, leftmargin=*}
\setlist[itemize]{itemsep=4pt, topsep=4pt, leftmargin=*}
\renewcommand{\arraystretch}{1.45}

% -- An expression the student is asked to work from -------------------------
\newcommand{\expr}[1]{%
  \tcbox[on line, colback=cbgreen!10, colframe=cbgreen!55, boxrule=1pt, arc=4pt,
         left=9pt, right=9pt, top=4pt, bottom=4pt]{\large\bfseries $#1$}}

% -- Running head and foot ---------------------------------------------------
% The packet overrides \fancyhead[L] with its own title.
\pagestyle{fancy}
\fancyhf{}
\fancyhead[L]{\small\color{black!55}Year 7\sep Homework}
\fancyhead[R]{\small\color{black!55}Name: \blank[3.4cm]}
\fancyfoot[C]{\small\color{black!55}Page \thepage}
\renewcommand{\headrulewidth}{0pt}
\renewcommand{\headrule}{\vspace{2pt}\hbox to\headwidth{\color{cbteal!45}\leaders\hrule height 1.2pt\hfill}}



% Landscape, with tighter margins than a packet: this page is a drawing surface.
% includeheadfoot is load-bearing -- without it geometry measures `top` to the
% top of the BODY and hangs the running head above it, off the top of the sheet.
\geometry{landscape, top=1.0cm, bottom=0.8cm, left=1.3cm, right=1.3cm,
          includeheadfoot, headsep=7pt}

\fancyhead[L]{\small\color{black!55}Algebra Track\sep 1.3 The Circle Map}
% The shared style puts a Name rule in the right of the running head. This sheet
% has one in its title block, and two on one page looks like a mistake.
\fancyhead[R]{}
\fancyfoot[C]{\scriptsize\color{black!40}Map: equidistant cylindrical blank map
  of Europe, public domain (Wikimedia Commons, user Anameofmyveryown)}

% -- The map overlay ---------------------------------------------------------
% \city{x}{y}{anchor}{name} -- x and y are fractions of the image box, anchor is
% the side of the DOT the label sits on (east = label to the left of the dot).
% The white fill behind the label is not decoration: without it a name landing
% on a coastline is unreadable, and half of these do (Copenhagen sits on the
% Danish islands, Stockholm on the Baltic, Athens in the Aegean).
\newcommand{\city}[4]{%
  \node[anchor=#3, inner sep=2.4pt, font=\sffamily\bfseries\footnotesize,
        text=black, fill=white, fill opacity=0.82, text opacity=1,
        rounded corners=1pt] at (#1,#2) {#4};
  \fill[cbred] (#1,#2) circle (1.7pt);}

\begin{document}
\thispagestyle{fancy}

% ===========================================================================
% TITLE
% ===========================================================================
\noindent
\begin{tcolorbox}[enhanced, nobeforeafter, width=\textwidth,
                  colback=cbpurple!7, colframe=cbpurple!7, arc=6pt,
                  borderline west={6pt}{0pt}{cbpurple},
                  left=16pt, right=12pt, top=6pt, bottom=6pt]
  {\color{cbpurple}\bfseries ALGEBRA TRACK\sep UNIT 1\sep LESSON 1.3}\quad
  {\color{cbpurple}\LARGE\bfseries The Circle Map}\hfill
  {\color{black!70}\bfseries Name: \blank[4.4cm]}
\end{tcolorbox}

\vspace{4pt}

% ===========================================================================
% MAP + INSTRUCTIONS
% ===========================================================================
% Both minipages open with \vspace{0pt} so [t] aligns their TOPS rather than the
% first baseline -- without it the tikzpicture's baseline sits at the foot of the
% image and the side column starts level with the bottom of the map.
%
% The widths have 0.5 cm of slack in them on purpose. City labels can stick a
% few points past the image edge, and at 18.6 cm + 7.5 cm the row came out
% 2.05 pt too wide, which silently wrapped the side column onto its own line and
% turned a one-page handout into three.
\noindent
\begin{minipage}[t]{17.2cm}
  \vspace{0pt}
  \begin{tikzpicture}
    \node[anchor=south west, inner sep=0, draw=rulegray, line width=0.5pt]
      % Reached across to the lesson unit rather than copied in beside this
      % file, so the sheet and the projector can never drift apart.
      (map) at (0,0) {\includegraphics[width=17.2cm]%
        {../../content/freshman-math/U01_3/images/europe-map.png}};
    \begin{scope}[x={(map.south east)}, y={(map.north west)}]
      \city{0.0681}{0.0973}{north}{Lisbon}
      \city{0.1975}{0.1577}{west}{Madrid}
      \city{0.3374}{0.1924}{west}{Barcelona}
      \city{0.5833}{0.2108}{west}{Rome}
      \city{0.8506}{0.0708}{east}{Athens}
      \city{0.3418}{0.4592}{east}{Paris}
      \city{0.6756}{0.4360}{west}{Vienna}
      \city{0.2827}{0.5538}{east}{London}
      \city{0.4025}{0.5845}{south}{Amsterdam}
      \city{0.6049}{0.5900}{west}{Berlin}
      \city{0.7860}{0.5796}{west}{Warsaw}
      \city{0.1367}{0.6196}{east}{Dublin}
      \city{0.5849}{0.7028}{west}{Copenhagen}
      \city{0.7160}{0.8332}{west}{Stockholm}
    \end{scope}
  \end{tikzpicture}
\end{minipage}%
\hfill
% The side column has to stay SHORTER THAN THE MAP (12.0 cm), or the whole row
% grows to fit it and the sheet spills onto a second page. At \small with the
% default hwbox padding these three boxes came out 15.6 cm; \footnotesize,
% tighter padding, small skips and fewer words bring them to about 11.5 cm.
\begin{minipage}[t]{8.2cm}
  \vspace{0pt}
  \tcbset{sidebox/.style={hwbox, fontupper=\footnotesize, breakable=false,
      left=7pt, right=7pt, top=9pt, bottom=5pt,
      before skip=0pt, after skip=5pt, colbacktitle=white}}

  \begin{tcolorbox}[sidebox, colback=cbpurple!6, colframe=cbpurple,
                    coltitle=cbpurple, title={Task 1 \textemdash{} by eye}]
    Pick any \textbf{three} cities that are not in a straight line, and write
    them in the table.

    Find the circle through all three: guess the centre, swing, adjust, swing
    again. \textbf{Three times, in three colours.}
  \end{tcolorbox}

  \begin{tcolorbox}[sidebox, colback=cborange!6, colframe=cborange,
                    coltitle=cborange, title={Task 2 \textemdash{} properly}]
    Now \key{construct} one instead of guessing.

    \textbf{1.} Perpendicular bisector of one pair.\\
    \textbf{2.} Perpendicular bisector of another pair.\\
    \textbf{3.} They cross at the \key{centre}.\\
    \textbf{4.} Needle there, pencil on a city, swing.

    \textbf{Leave every arc on the page.}
  \end{tcolorbox}

  \begin{tcolorbox}[sidebox, colback=cbgreen!7, colframe=cbgreen,
                    coltitle=cbgreen!75!black, title={Challenge}]
    Find three cities you \textbf{cannot} draw a circle through. Name them, and
    say why in one sentence.

    \vspace{1pt}
    \blank[7.1cm]

    \vspace{3pt}
    \blank[7.1cm]
  \end{tcolorbox}
\end{minipage}

\vspace{6pt}

% ===========================================================================
% RECORDING TABLE
% ===========================================================================
% The shared style sets \arraystretch to 1.45, which is right for a Year 7
% answer table and wrong here: at four rows it cost 1.4 cm and pushed the sheet
% onto a second page. Row height is set explicitly with \rule instead, so the
% writing space is chosen rather than inherited.
\noindent
\renewcommand{\arraystretch}{1.0}
% tabularx sizes X to make the CONTENT total \textwidth and then adds the five
% vertical rules on top, so a plain \textwidth here overfills by 2.05 pt every
% time. Take the rules off the target width instead.
\begin{tabularx}{\dimexpr\textwidth-5\arrayrulewidth\relax}{|p{1.7cm}|X|X|p{3.0cm}|}
  \hline
  \rule[-0.4em]{0pt}{1.6em}\textbf{Circle} & \textbf{The three cities} &
  \textbf{The centre landed in / near} & \textbf{Radius (cm)} \\
  \hline
  \rule[-0.55em]{0pt}{1.75em}1 & & & \\
  \hline
  \rule[-0.55em]{0pt}{1.75em}2 & & & \\
  \hline
  \rule[-0.55em]{0pt}{1.75em}3 & & & \\
  \hline
\end{tabularx}

\end{document}
